\newcommand{\figuraTanqueTuboCapilar}{%
\begin{figure}[H]
\centering
\begin{tikzpicture}[>=Latex,font=\small,line cap=round,line join=round]
  \draw[very thick] (1.0,0.6) rectangle (7.2,4.5);
  \fill[blue!12] (1.08,0.68) rectangle (7.12,2.55);
  \draw[very thick,blue!65!black] (1.0,2.55)--(7.2,2.55);
  \node at (4.1,3.55) {$p_1$};
  \draw[very thick] (7.2,0.88)--(8.3,0.88)--(8.3,5.55);
  \draw[very thick] (7.2,1.22)--(8.05,1.22);
  \draw[very thick] (8.75,1.05)--(8.75,5.55);
  \fill[blue!18] (7.2,0.96) rectangle (8.3,1.14);
  \fill[blue!18] (8.38,0.96) rectangle (8.67,4.45);
  \draw[very thick,blue!70!black] (8.38,4.45)
    .. controls (8.45,4.32) and (8.60,4.32) .. (8.67,4.45);
  \draw[<->,red!75!black,thick] (9.35,2.55)--(9.35,4.45)
    node[midway,right] {$15\ \mathrm{cm}$};
  \draw[<->,thick] (8.3,5.85)--(8.75,5.85)
    node[midway,above] {$1\ \mathrm{mm}$};
  \draw[densely dashed] (7.0,2.55)--(9.6,2.55);
  \node[align=left,anchor=west] at (9.85,4.45)
    {nivel medido\\[1mm](incluye capilaridad)};
\end{tikzpicture}
\caption{Tanque conectado a un tubo piezométrico capilar. La lectura de
$15\,\mathrm{cm}$ combina la altura debida a la presión del tanque con la
corrección capilar del tubo de $1\,\mathrm{mm}$.}
\label{fig:tanque-tubo-capilar}
\end{figure}%
}